\subsection{The two remaining coned three-state star cases}

Write $P=P_5=(\{0,1\}\times\{0,1\})^1$, with
$e_{ij}e_{kl}=e_{il}$.  For a left-$N$, right-$M$ biset $S$,
put
\[
 E=B(N,S,M),\qquad Y=\operatorname{hocofib}(E\longrightarrow BN).
\]
This is the classifying-space model of the source star obtained
by adjoining to the $(M,N,S)$ edge a trivial leaf on a singleton
stalk: that additional edge cones off $BM$.
All homology in this subsection is rational unless a coefficient
ring is displayed.

\begin{lemma}[Proved: elementary coned-bridge comparisons]
\label{CP2:A-coned-bridge-comparisons}
The following assertions hold for the preceding model.
\begin{enumerate}
\item If a second class of $Y$ restricts to zero on $BN$, its
square is zero.  In particular all second squares vanish if
$H^2(BN)=0$, or if $E\to BN$ has a section on second cohomology.
\item If the left $N$ action on $S$ is trivial, then $Y$ is
homotopy equivalent to a suspension.
\item A proper biset retraction of $S$ can be used in a
projective-plane exclusion: a star of type $\CP^2$ would give
a star of the same type on the image biset.  In particular a
retraction onto at most two states contradicts
Lemma~\ref{CP2:A-star-coned-binary}.
\item Suppose $BM$ is rationally acyclic, the right action
space $Z=B(S,M)$ is connected, and some left $N$ actor on $S$
is constant.  Then $H_2(Y)=0$.
\end{enumerate}
\end{lemma}
\begin{proof}
The quotient $Y/BN$ is the suspension of $E$. A class with
zero restriction lifts from this quotient, whose positive
products vanish. The cofiber exact sequence gives the second
assertion in (1).

For (2), $E=BN\times Z$ and its map to $BN$ is projection.
Choose a vertex of $Z$. Restricting the projection to the
resulting section gives the identity $BN\to BN$; its mapping
cone is a contractible subcomplex of the full cofiber model.
Collapsing it identifies the quotient with
$\Sigma(BN_+\wedge Z)$, with the chosen vertex used as the
basepoint of $Z$. Collapsing that contractible subcomplex is
a homotopy equivalence.

For (3), use Theorem~\ref{CP2:A-cut-retract} for the cut
from the center to its two leaves, with the identity on the
singleton stalk. This statement concerns an actual retract,
so the integral projective-plane recognition in that theorem
applies.

For (4), regard $Z$ as a left $N$ space. The cone of the
augmented complex $C_*(Z)\to\mathbb Q$ has no homology in
degrees zero and one, and its degree-two homology is $H_1(Z)$.
The derived $N$ coinvariants of this cone compute the reduced
homology of $Y$. The first nonzero hyperhomology row therefore
gives $H_2(Y)=H_1(Z)_N$. A constant actor with value $u$ has
$u m=u$ for every $m\in M$, by commutation. Its map on $Z$
factors through $BM$, followed by the fixed-state inclusion
$BM\to Z$. Thus it acts by zero on $H_1(Z)$. Its coinvariant
relation kills every element, proving (4).
\end{proof}

\begin{proposition}[Proved: a P5 target on at most three states]
\label{CP2:A-coned-three-target-P5}
If $N=P$ and $1\leq |S|\leq3$, the coned bridge $Y$ cannot
have the homotopy type of $\CP^2$. There is no restriction on
the finite center monoid $M$.
\end{proposition}
\begin{proof}
A state fixed by the left $P$ action gives a map $BP\to E$
using empty $M$ strings. Its composite with the projection
to $BP$ is the identity; the state need not be fixed by $M$.
Lemma~\ref{CP2:A-coned-bridge-comparisons}(1) then applies.

Otherwise the left rectangular-band act normal form has
a unique two-point regular row core $U=\{c_0,c_1\}$ and at
most one transient point. The left band acts by
$e_{ij}s=c_i$. Every commuting right $M$ actor fixes $U$
pointwise. Swapping the two band columns defines an
automorphism $F$ of the star that fixes $M$ and $S$.

Use the opposite of the length-two projective resolution
in Lemma~\ref{CP2:A-order-five-corner}, and tensor it with
$\mathbb Q[S]$. Quotienting the acyclic subcomplex supported
on the two core points gives the right $\mathbb QM$ complex
\[
 \mathbb Q[0]\oplus T[1],\qquad T=\mathbb Q[S]/\mathbb Q[U],
\]
with zero differential. The column-interchange semilinear
lift, the opposite of the lift displayed in
Lemma~\ref{CP2:A-mixed-five-reductions}, acts as $+1$ on
$\mathbb Q[0]$ and as $-1$ on $T[1]$. Derived tensoring with
the trivial left $M$ module consequently gives the natural
decomposition
\[
 H_k(E)=H_k(BM)\ \oplus\
 \operatorname{Tor}^{\mathbb QM}_{k-1}(T,\mathbb Q),
\]
with respective $F$ signs $+1$ and $-1$. The second term is
absent in degree zero. The core bridge $E_U=B(P,U,M)$ maps
onto exactly the first summand.

Here the geometric core comparison is needed as well.
The left act $U$ is $Pe_{00}$, a retract of the regular left
act $P$. Its action category is therefore contractible: the
regular action category is contractible, and the act inclusion
and retraction give an actual categorical retract. As $M$
fixes every core point,
\[
 E_U=B(P,U)\times BM\simeq BM.
\]
Its projection to $BP$ factors through the contractible
first factor and is nullhomotopic. This nullhomotopy yields
a map $\Sigma E_U\to Y$ whose connecting homomorphism to
$H_*(E)$ is the core inclusion. Since $H_3(BP)=H_4(BP)=0$,
the cofiber sequence identifies $H_4(Y)$ with $H_3(E)$.
Thus this suspension map realizes the entire $+1$
eigenspace $H_3(BM)$ in $H_4(Y)$.

Suppose $Y\simeq\CP^2$ and let $\alpha$ generate $H^2(Y)$.
If its restriction to $BP$ vanishes, its square is zero by
Lemma~\ref{CP2:A-coned-bridge-comparisons}(1). Otherwise the
nonzero restriction into the one-dimensional $H^2(BP)$,
on which column interchange is $-1$, implies
$F^*\alpha=-\alpha$. Hence $F^*(\alpha^2)=\alpha^2$.
Its pairing with the $-1$ eigenspace of $H_4(Y)$ is zero.
Its pairing with the $+1$ eigenspace is also zero, because
that eigenspace is realized by $\Sigma E_U$, where every
second square vanishes. Over $\mathbb Q$ the two eigenspaces
span $H_4(Y)$; therefore $\alpha^2=0$, a contradiction.
\end{proof}

\begin{proposition}[Proved: a P5 center on at most three states]
\label{CP2:A-coned-three-center-P5}
If $M=P$ and $1\leq |S|\leq3$, the coned bridge $Y$ cannot
have the homotopy type of $\CP^2$. There is no restriction on
the finite leaf monoid $N$.
\end{proposition}
\begin{proof}
The right $P$ fixed states form an $N$ invariant subset.
If it has a single element, that element is also $N$ fixed
and gives a section $BN\to E$, so all second squares vanish.

Suppose there are at least two fixed states. Cardinality
three leaves no two-point regular core orbit in addition,
so the entire band core is the fixed set $F$. Put
$X=\mathbb Q[S]$ and $U=\mathbb Q[F]$. Tensoring the explicit
left $\mathbb QP$ resolution with $X$ gives the $\mathbb QN$
complex
\[
 U\oplus U \xrightarrow{(v,w)\mapsto v+w} X
 \xrightarrow{d_1} U,
 \qquad d_1(s)=s e_{10}-s e_{00},\qquad d_1|_U=0,
\]
in degrees two, one and zero. Quotient its acyclic subcomplex
given by the first $U$ in degree two mapping identically onto
$U\subset X$. The quotient is
\[
 U[2]\ \oplus\ \bigl[X/U\xrightarrow{\bar d_1}U\bigr],
\]
where the bracket occupies degrees one and zero. The bridge
projection to $BN$ uses only the degree-zero augmentation;
it is zero on the displayed $U[2]$ summand. After derived
$N$ coinvariants and passage to the cofiber, $H_3(Y)$ has
the direct summand $U_N$. This is nonzero because $F$ is a
nonempty finite $N$ set. Thus $Y$ does not have the homology
of $\CP^2$.

It remains that there are no fixed states. There is then
one regular two-point column core $U=\{c_0,c_1\}$ and at
most one transient point. Every $N$ actor fixes the core
pointwise. The acyclic-core quotient of the same projective
resolution is
\[
 \mathbb Q[0]\oplus T[1],\qquad
 T=\mathbb Q[S]/\mathbb Q[U].
\]
Under this quotient the projection $E\to BN$ is exactly the
projection onto the first summand. Therefore
$H_2(Y)=\mathbb Q\otimes_{\mathbb QN}T$. If $T=0$ this group
is zero. Otherwise $T$ is one-dimensional: an $N$ actor has
scalar one if it fixes the unique transient state and scalar
zero if it sends that state into the core. Nonzero coinvariants
force every scalar to be one. All states are then $N$ fixed,
so Lemma~\ref{CP2:A-coned-bridge-comparisons}(2) makes $Y$ a
suspension. This excludes the final case.
\end{proof}

\begin{proposition}[Proved: the other five-element centers]
\label{CP2:A-coned-three-other-five}
Let $|M|=5$, $M\not\cong P$, $|N|\leq11$, and
$1\leq |S|\leq3$. The coned bridge $Y$ cannot have type
$\CP^2$.
\end{proposition}
\begin{proof}
The center $BM$ is rationally acyclic by
Lemma~\ref{CP2:A-order-five-corner}. If its minimal ideal
has one row, the trivial left $\mathbb QM$ module is
projective, by normalized local-group averaging. The actual
right-orbit quotient $S\to S/M$ induces a rational homology
and cohomology-ring equivalence from $Y$ to the ordinary
left $N$ act cone on $S/M$. If $Y$ were a projective plane,
this latter cone would have bounded rational homology.
The left-act version of Theorem~\ref{CP2:A-bounded-eleven-acts}
kills every second product there. Pullback gives a
contradiction. This comparison requires only rational
boundedness, not integral projective-plane recognition.

Otherwise the minimal ideal has at least two rows. Since
$M\ne P$ and $|M|=5$, it has exactly one column. First suppose
its local group is trivial. Every core state is then globally
$M$ fixed: right multiplication fixes every element of a
one-column rectangular band. Conversely every fixed state
belongs to the core. Write $F$ for this nonempty fixed set.
If $|F|=1$, commutation makes its point $N$ fixed and gives
the section in Lemma~\ref{CP2:A-coned-bridge-comparisons}(1).
If $F=S$, the $M$ action is trivial. The bridge is
$B(N,S)\times BM$, and rational acyclicity of $BM$ again
gives a rational ring equivalence to the ordinary $N$ act
cone on $S$, excluded by bounded-eleven.

Only $S=\{u_0,u_1,w\}$ with $F=\{u_0,u_1\}$ remains.
The right transformation image fixes both core points and
sends $w$ to $w,u_0$, or $u_1$. If it has only one
nonidentity retraction $p$, then $p$ commutes with the whole
right image and with every $N$ actor. It is a biset retraction
onto two states, so the coned-binary lemma and the cut-retract
theorem exclude an actual projective-plane type.
Otherwise both retractions $p_0(w)=u_0$, $p_1(w)=u_1$ occur.
A map commuting with both is either the identity or a full
constant to $u_0$ or $u_1$: its value on $w$ determines its
values on both core points. If the $N$ image is the identity
alone, $Y$ is a suspension. If that image contains a constant,
the right action space is connected because $w$ reaches both
core points. Lemma~\ref{CP2:A-coned-bridge-comparisons}(4)
then gives $H_2(Y)=0$.

Finally suppose the one-column ideal has a nontrivial local
group. Its size and the outside identity force
$M=(\{0,1\}\times C_2)^1$, with
$(i,g)(j,h)=(i,gh)$. Its core is a disjoint union of $C_2$
orbits of sizes one or two. A single globally fixed state is
$N$ fixed, and the preceding section argument applies. If
all core orbits have one point, the transformation analysis
of the preceding paragraph applies verbatim, since the group
acts trivially on the core. The remaining possibility has a
single regular two-point orbit and at most one transient $w$.
With no transient the common core retraction already has two
states. With a transient, let $p_0,p_1$ be the two ideal
idempotent actors. They retract onto the same two-point core.
If $p_0(w)=p_1(w)$, their common retraction commutes with the
core transposition and all ideal actors, hence with $M$ and
$N$; apply coned-binary. If their values differ, a commuting
map sending $w$ into the core would identify the two core
points, contradicting commutation with the core transposition.
A commuting map fixing $w$ must fix both core points. Thus
the entire centralizer is the identity and the $N$ action is
trivial. Its cofiber is a suspension, completing the proof.
\end{proof}

\begin{corollary}[Proved: both outstanding large star profiles]
\label{CP2:A-star-coned-three-five-eight}
The source-star profiles
\[
 (\text{center},\text{leaf}_1,\text{leaf}_2,
   \text{stalk}_1,\text{stalk}_2)
 =(5,1,8,1,3),\qquad(8,1,5,1,3)
\]
and their leaf interchanges have no classifying space of type
$\CP^2$.
\end{corollary}
\begin{proof}
The singleton trivial leaf cones the center, giving precisely
the model $Y$. In the first profile use
Propositions~\ref{CP2:A-coned-three-center-P5} and
\ref{CP2:A-coned-three-other-five}. In the second profile a
non-$P$ five-element leaf has zero second rational cohomology,
so Lemma~\ref{CP2:A-coned-bridge-comparisons}(1) applies;
the remaining leaf $P$ is Proposition~\ref{CP2:A-coned-three-target-P5}.
\end{proof}

\begin{remark*}[Correction of a discarded route]
No graph model is asserted for a $P$ act with two globally
fixed core states and one transient state. In particular, when
the two row retractions send that transient to different core
states, the displayed projective complex has $H_0=\mathbb Q$,
$H_1=0$ and $H_2=\mathbb Q^2$; it is not an interval. An
initial proposed graph shortcut was withdrawn before use.
The degree-two module summand in the center-$P$ proof is the
replacement argument, and the regular-core proof for a target
$P$ uses the actual contractible left action category only.
\end{remark*}
